\subsection{Voting Data Integration}
\noindent The translation of neighborhood protest into formal City Council voting behavior is tracked via longitudinal voting integration. Voting data extracted from the Council Meeting System (CMS) confirms that the degradation of council consensus is strongly associated with the presence of formal protest petitions. Unprotested discretionary zoning changes are typically passed on consent (unanimously), whereas petitioned cases frequently fracture the council vote. Table~\ref{tbl:voting_data} aggregates these voting outcomes across the longitudinal panel.

\begin{table}[h!]
\centering
\begin{tabular}{l c c}
\toprule
\textbf{Voting Context} & \textbf{Unprotested Cases} & \textbf{Petitioned Cases} \\
\midrule
Passed on Consent (Unanimous) & 92\% & 31\% \\
Contested Vote ($<$100\% agreement) & 8\% & 69\% \\
\bottomrule
\end{tabular}
\caption[Council Voting Consensus by Protest Status]{Council Voting Consensus by Protest Status. Sourced from the integrated CMS voting panel, illustrating the collapse of unanimous consent when a valid petition is filed.}
\label{tbl:voting_data}
\end{table}
